\section{Changing carriers relative to input-dependent references}\label{sec:references}

Public catalogues need not contain all raw-dependent words, even when each
word has a simple local description. We therefore use a counted analysis
model with one physical work tape, $p$ continuous raw-window ports, the
strict prior output, and one immutable input-dependent label. Its whole
transition may use these arguments. Its finite control has $2^{q_*}$ states;
time, positions and the append ordinal are explicit arguments, and every
other evolving register is charged. It starts in a public state with blank
work tape and public port origins. Every work write is at its old position,
all heads and ports move by at most one, and at most one bit is appended.
This model is only an intermediate counting device.

\begin{lemma}[Counting the entire enlarged transition]\label{lem:multiport-transition}
\formal{EnlargedFinite.lean}{Enlarged.finite_bound}
The uniform tail is \lean{Enlarged.macroscopic_tail} in
\texttt{EnlargedTail.lean}.
Let $P=d^2b\ge4p$, $T=\lfloor C(P+1)\rfloor$,
$N=\lfloor640apC^2\rfloor+1$, $W=\lceil P/(320apC)\rceil$, and
$R\le\lfloor P/(8pN)\rfloor$. A fixed system of the preceding kind with
alphabet $\F_2^a$ and label support at most $2^z$ correctly transposes, for
some label, at most
\[
 2^{P-P/(4N)+3Ndb+a+Nq_*+z+\log_2\mathcal M_1}
\]
inputs, where
\[
 \mathcal M_1=(N+1)(2T+W+2)(T+1)^2
       [(T+1)^4(2T+1)^{1+3p}]^N.
\]
If $q_*\le P/(64N^2)$ and $z\le P/(64N)$, this is less than
$2^{(1-1/(8N))P}$ uniformly in $b\ge1$ at sufficiently large $d$.
\end{lemma}
\begin{outline}
Apply the slab selection and raw-footprint calculation in
\cref{thm:getter-tail} to the physical work head. The selected chunk
has at most $N$ runs and saves $P/(4N)-3Ndb-a$ bits after its complete
initial slab is paid. There are no noncarrier symbols or error tags.
Record run-entry control, port positions/extrema, physical work position,
times and append counters, and one label. This is $\mathcal M_1$.
Replay the entire transition from the persistent work slab and the known
raw windows. Assemble its strict prior output exactly as in
\cref{thm:getter-replay}; this occurs before applying the transition, so
no current output bit is assumed. The same induction applies to the
whole enlarged rule, including its port updates. Correctness recovers the
omitted input. The label is paid once and run-entry controls at most $N$
times, giving the stated finite bound. The two linear budgets and
$3Ndb+a+\log_2\mathcal M_1=o(P)$ leave the claimed slack uniformly in $b$.
\end{outline}

Fix $k\ge2$ native work tapes, $\ell\ge1$ references per tape, and an
embedding of the source alphabet in $V=\F_2^a$ with blank zero. Put
$E=[-T,T]\cap\Z$, $n=2T+1$. Each reference uses a public $1$-Lipschitz
map $\phi_{hi}:\Z\to\Z$ and a public local function $\gamma_{hi}$ of
the coordinate and the radius-$r_0$ raw window centered at $\phi_{hi}(z)$.
One optional immutable seed $\lambda_0=f(X)$ is fixed before the hash used
below. Its joint support on the fixed promise set is at most $2^{Z_0}$.
An observer with $U$ states contains all cutoffs $b_{hi}\in[0,P]$ and
auxiliary storage, starts publicly, and changes each cutoff by at most one.
Its updates may use the seed, old source/observer context, strict output,
actual old source symbols and the declared old raw windows, but no other
individual raw data. Initial cutoffs and port origins are public, independent
of the seed. All rules are independent of the later hash choice.

Let $D_{hi}(t,Y,\lambda_0)\in V^E$ be any public dense mask. At reached
time $t$ define
\[
 F_{hi}(t,X)(z)=\mathbf1_{[0,b_{hi}(t))}(z)
   \gamma_{hi}(z,\operatorname{window}_X(\phi_{hi}(z)),\lambda_0)
      +D_{hi}(t,Y_t,\lambda_0)(z),
\]
where $Y_t$ is strictly prior output. Words are zero outside $E$.
Write
\[
 K_s=\sum_{j=0}^{\min(s,n)}\binom nj(2^a-1)^j,
 \quad r=\lceil\log_2(2k(T+1)\ell K_s)\rceil,
 \quad p=1+2k\ell.
\]

\begin{theorem}[Raw-relative carrier handoff]\label{thm:reference-carrier}
\notformalized
Let $G$ be any promise set of source inputs whose runs halt by $T$, append
at most $P$ bits, obey the observer promises, and at every reached
configuration, including halt, have at most one work word farther than
$s$ cells from every current reference on its tape. One common analysis
system of \cref{lem:multiport-transition}, with $p$ ports and radius
$r_0+1$, reproduces the source exactly on at least $|G|/2$ inputs, with
the same transition count, raw input trajectory and append times/bits.
Its state and label bills are
\[
 Q_*\le2+QU B_{\rm ref}2^{k(\ell+1)r},\qquad
 z\le Z_0+k\ell r,
\]
where $Q$ is the number of source states and
\[
 B_{\rm ref}=k(T+1)n^{k+2}(4T+1)(1+n^2)
             (1+T(P+1)^2\ell K_s n2^a).
\]
No ordinary native one-tape simulation is asserted.
\end{theorem}
\begin{informalproof}
Choose a binary linear hash $H:V^E\to\F_2^r$. For fixed $X\in G$, its
original augmented run and seed are determined independently of $H$.
Compare each actual word $w_h(t)$ with every different candidate
$F_{hi}(t,X)+e$, where $e$ has at most $s$ nonzero cells. Each nonzero
difference hashes to zero with probability $2^{-r}$: in each independently
chosen binary row of $H$, one nonzero coordinate determines a fair bit.
There are at most $k(T+1)\ell K_s$ comparisons, so a union bound followed
by averaging supplies one $H$ good on at least half of $G$.
This event refers only to original trajectories. Empty $G$ is immediate.
The same finite selection theorem \lean{exists_binary_common_half} applies
with $B=k(T+1)\ell K_s$. It proves the single-hash half-set claim, including
empty comparison families; the native reference-system construction below
remains unformalized.

Maintain $k$ actual-word hashes $Z_h=H(w_h)$ and $k\ell$ raw-prefix
hashes $L_{hi}=H(\mathbf1_{[0,b_{hi})}\gamma_{hi})$. Supply their initial
raw values by one label together with $\lambda_0$; this costs at most
$Z_0+k\ell r$, and their evolving copies are separately in the controller.
The actual hashes initially vanish. An old-symbol change $u\to v$ at $y$
updates $Z_h$ by $H(\mathbf1_{\{y\}}(u+v))$, even for dense words.
A cutoff increment adds the hash contribution at its old boundary $b$;
a decrement removes the contribution at $b-1$. Contributions outside $E$
are zero. Use continuous ports for the source input, every $\phi_{hi}(y_h)$
and every $\phi_{hi}(b_{hi})$. Their unit motion follows from the
$1$-Lipschitz maps. The old radius-$(r_0+1)$ boundary window contains the
decrement's required radius-$r_0$ window. No port jumps at handoff.

Test, in a public fixed order, the finite equations
\[
 H(e)=Z_h+L_{hi}+H(D_{hi}(t,Y_t,\lambda_0)).
\]
On a good original run every matching pair describes the actual word and
an unmatched word has none. Several pairs describing the same word are
allowed; choose the first. Candidate words outside the source alphabet may
be included in this test, so no global raw-word legality query is needed.
Only the reconstructed scanned symbols need belong to the source alphabet.

Keep the active index $f$, logical source positions, physical position $p_0$,
offset $u$ with $y_f=p_0+u$, and an empty or integer interval $I$ of physical
positions scanned in previous transitions of this carrier phase. Its frozen
entry word is initially blank, or is encoded by
\[
 (e,o_e,b,i,e_s,z,v),\quad e<T,\quad o_e,b\in[0,P],
 \quad e_s\in\operatorname{Err}_s,
\]
meaning the old reference with cutoff $b$, mask
$D_{fi}(e,Y_t[0:o_e),\lambda_0)$ and error $e_s$, with cell $z$
\emph{overwritten} by $v$. For $j\in I$ the physical symbol is the current
logical symbol at $j+u$; outside $I$ the logical word still equals this
frozen entry word. Its fresh symbol uses the logical-head raw port and
the historical prefix slice, never the current cutoff or mask.

In a public INIT state install the public source/observer data, initialize
hash registers from the label and choose the unique unmatched tape, or any
tape if all match. This initialization is part of evaluating the first
ordinary transition, with no extra step or free tape write. Initially all
logical words are blank, $u=0$ and $I$ is empty. Invalid data can halt.

At a represented step decode every inactive scanned symbol from a matching
pair. Decode the active symbol physically when $p_0\in I$, and from its
frozen code otherwise. Having the entire old tuple, apply the source rule,
observer update and the local hash updates. A source halt causes no write
or append. Otherwise the sole physical action writes the old active tape's
new symbol at $p_0$ and follows that tape's chosen move. All logical positions
and raw ports make their unit moves, and the same bit is appended at the
same step. Only now test the post-step references using the new prefix.

If the active tape stays selected, extend $I$ by the old physical position;
unit motion keeps it an interval. If a new carrier $g$ is selected, it was
inactive before this transition. Freeze its PRE matching pair, old cutoff,
time and output ordinal, together with its one actual old-position write.
This describes its exact POST word. Set the new offset from its new logical
position and the already moved physical head, and reset only $I$ to empty.
The former carrier is now matched. Stale physical cells are ignored on first
visit and overwritten with their correct next values; no clearing or seek
is performed. These statements establish the invariant by induction through
arbitrary switches, including changes caused solely by masks or cutoffs.

Malformed fields, nonunit moves or unavailable matches can halt; no branch
tests membership in the good set. The invariant proves exactness on that
set with its true label. Count source/observer state, active index, time,
$k+2$ positions in $[-T,T]$, offset in $[-2T,2T]$, interval, the
$k(\ell+1)$ hashes and the frozen code. Their product is the displayed
bound, with INIT and a rejecting state added. Port positions are derived
from counted logical positions/observer state; the strict prefix itself is
an analysis argument. The immutable seed support can be counted on the fixed
$G$ before choosing $H$, with no second largest-fibre loss.
\end{informalproof}

\begin{corollary}[Two words far from all current references]\label{cor:reference-tail}
\partformal{ReferenceTail.lean}{reference_tail_of_half_simulation,reference_cache_sparse_reserves}{Lean derives the complete state and label budgets, output correctness and success tail from one common native-prefix simulation on at least half the promise set, with the literal cardinality bills. The reference-carrier construction supplying that simulation remains unformalized.}
Fix the above resources, set $N=\lfloor640apC^2\rfloor+1$ and
$D_0=k(\ell+1)+1$. If
\[
 q_{\rm src}+q_{\rm obs}\le\frac P{128N^2},\quad
 \lceil\log_2 K_s\rceil\le\frac P{256D_0N^2},\quad
 Z_0\le\frac P{128N},\quad r_0+1\le\left\lfloor\frac P{8pN}\right\rfloor,
\]
then fewer than $2^{1+(1-1/(8N))P}$ timely correct inputs satisfy the
one-exception promise, uniformly for all $b\ge1$ at large $d$.
For a sufficiently small fixed $\delta>0$, $s=\lfloor\delta P\rfloor$
meets the sparse budget. Outside that exception, a reached configuration
has two distinct words each farther than $s$ cells from every current
reference for its tape.
\end{corollary}
\begin{informalproof}
Put $\lambda_s=\lceil\log_2K_s\rceil$ and
$L_0=\lceil\log_2(2k(T+1)\ell)\rceil$. Then $r\le\lambda_s+L_0$,
and $B_{\rm ref}\le K_s B_0$ for the same polynomial factor with
$K_s$ removed. The compiler therefore has
\[
 q_*\le q_{\rm src}+q_{\rm obs}+D_0\lambda_s+O(\log(P+1))
       \le P/(64N^2)
\]
eventually, uniformly in $b$. The same estimates give
$k\ell r\le P/(128N^2)$, hence $z\le P/(64N)$ including $Z_0$.
Apply \cref{lem:multiport-transition} and retain the factor two from the
common-hash half-set. For the last assertion use
$n\le\lceil4C+1\rceil P$ and
$\log_2K_s\le s\log_2(en(2^a-1)/s)$ for sufficiently small $s/n$.
Its ratio to $P$ tends to zero with $s/P$; choose a fixed positive
$\delta$ with enough slack for the ceiling. Negating the one-exception
promise gives the stated simultaneous witness. This is per declared system,
not a union over input-selected reference rules.
The checked budget includes both administrative states and every integer
ceiling. It uses a single dimension threshold before the width, source,
observer and enlarged transition table are chosen.
\end{informalproof}

\subsection{Static scans as input-dependent references}

Take $L=\lfloor P/(8pN)\rfloor$. The checkpoint construction of
\cref{lem:scan-checkpoints} with gaps at most $L$ evaluates each static
original-raw scan using radius $L-1$; cutoff decrements then use total radius
$L$. Its one joint seed is reused at every handoff. It fits the preceding
corollary when its support charge is at most $P/(128N)$. In particular,
for $P\ge(8pN)^2$ the number of equal-grid internal anchors is
$\lfloor(P-1)/L\rfloor\le8pN$: writing $P=8pN L+r$, $0\le r<8pN$
and $L\ge8pN$ proves the inequality. Thus common state bits
$w\le P/(1024pN^2)$ suffice. Public offset averaging also gives the
sufficient profile condition $\sum_i w_i\le LP/(128N)$.

The scan descriptions remain functions of the original raw input; the
current cutoffs, masks and sparse deviations may evolve. A large offline
scan table has no native production guarantee. Public catalogues with a
growing number of references retain their separate scope, as do fixed-carrier
local discrepancies. Neither conclusion supplies independence, persistence
or output usefulness of the deviations. A raw-copy track alongside an
arbitrary payload does not make that payload a small residual.
